\section{A bounded flat arity block cannot retain native torsion}
\label{sec:136-flat-arity-tail}

\paragraph{Archive scope.}
The application concerns the vertical uniformity condition on cote 136,
physical p.~117, footer 116, supplied scan 117, and the common-source
question on physical p.~121, footer 120, scan 121. It tests a specific
proposal: replace the strict divisibility relations of the radial disk
envelope by a free homogeneous resolution, retaining the auxiliary
divided-power arity. It does not exclude arbitrary nonfree completions,
constructions mixing arities, or constructions unbounded in form degree
at a fixed arity.

\paragraph{ARCH-136-FLAT-TAIL-ACYCLICITY.}
Let \(k=\mathbb F_p\), and let
\(\overline S=\Gamma_k(T)\), with \(T\) of weight one.
Let \(F\) be a bounded cochain complex of weight-graded flat
\(\overline S\)-modules, with weight-preserving differential.
Suppose there are finite integers \(L\leq U\) such that
\[
                 H^j(F)^w=0\quad\text{if }w\notin[L,U]
                       \quad\text{for every }j.
                                                        \tag{136.FT1}
\]
Then \(F\) is acyclic. Ranks of modules and dimensions of the
cohomology groups in this finite band may be arbitrary.

\begin{proof}
Choose \(a\geq0\) such that \(h=p^a>U-L\), and set
\[
 x=T^{[p^a]},\qquad D=k[x]/(x^p),\qquad\operatorname{wt}(x)=h.
                                                        \tag{136.FT2}
\]
There is a weight-graded decomposition
\[
 \overline S\cong\bigotimes_{b\geq0}
          k[x_b]/(x_b^p),\qquad x_b=T^{[p^b]}.
                                                        \tag{136.FT3}
\]
Indeed, writing \(w=\sum e_bp^b\), \(0\leq e_b<p\), the
monomial \(\prod_b x_b^{e_b}\) is a nonzero scalar multiple of
\(T^{[w]}\): the valuation of its integer coefficient is zero
by the factorial valuation formula. The coefficient of \(x_b^p\)
has valuation one. These assertions give one basis vector in each
weight and prove (136.FT3). Thus \(\overline S\) is free over
the selected factor \(D\), and each \(F^j\) is flat over \(D\).

Multiplication by \(x\) is zero on every \(H^j(F)\), because it
raises weight by more than the entire band (136.FT1). The standard
graded free resolution \(P_\bullet\to k\) over \(D\) has alternating
maps \(x,x^{p-1}\). Its generator weight shifts are
\[
 \sigma_{2r}=rph,\qquad \sigma_{2r+1}=rph+h
                    \quad(r\geq0).                    \tag{136.FT4}
\]
The resolution is exact since the annihilators of \(x\) and
\(x^{p-1}\) are respectively \((x^{p-1})\) and \((x)\).
All its maps induce zero after tensoring with \(H^j(F)\).

Consider the double complex \(P_i\otimes_D F^j\), with total
cochain degree \(j-i\). Since \(F\) is bounded, every total
degree has a finite diagonal. Computing its cohomology first in
the \(F\)-direction gives the hyper-Tor spectral sequence
\[
 E_2^{-i,j}=\operatorname{Tor}_i^D(k,H^j(F))
   \ \Longrightarrow\ H^{j-i}(k\otimes_D F).
                                                        \tag{136.FT5}
\]
The displayed abutment can also be checked directly: computing first
in the resolution direction gives \(k\otimes_D F\), because each
\(F^j\) is \(D\)-flat. No unbounded product totalization is involved.

The \(i\)-th column of (136.FT5) is supported in the weight band
\([L+\sigma_i,U+\sigma_i]\). The shifts are strictly increasing,
with consecutive differences at least \(h>U-L\). These bands are
pairwise disjoint. Every spectral-sequence differential changes the
column and preserves weight. Consequently every differential is zero,
and \(E_2=E_\infty\).

If \(H^j(F)\neq0\), its copy occurs in every column \(i\geq0\).
It would therefore give nonzero abutment groups in the arbitrarily
negative cochain degrees \(j-i\). This is impossible, because
\(k\otimes_D F\) is a bounded cochain complex. Hence all
\(H^j(F)\) vanish, as asserted.
\end{proof}

\paragraph{ARCH-136-INTEGRAL-FLAT-TAIL.}
Let \(R=\mathbb Z_{(p)}\), \(S_R=\Gamma_R(T)\), and let \(K\)
be a bounded complex of weight-graded flat \(S_R\)-modules.
Suppose its cohomology has a common finite weight band as in
(136.FT1). Then multiplication by \(p\) is an isomorphism on
every \(H^j(K)\). In particular these groups contain no nonzero
\(p\)-power torsion. If \(K\otimes_R\mathbb Q\) is acyclic,
then \(K\) itself is acyclic.

\begin{proof}
The terms of \(K\) are \(R\)-flat, so
\[
 0\longrightarrow K\xrightarrow{p}K
              \longrightarrow\overline K\longrightarrow0,
 \qquad\overline K=K\otimes_R\mathbb F_p,
                                                        \tag{136.FT6}
\]
is exact as a sequence of complexes. The coefficient long exact
sequence shows that the cohomology of \(\overline K\) has the
same common finite weight band. Its terms are flat over
\(S_R/pS_R=\Gamma_{\mathbb F_p}(T)\), by base change of flat
modules. ARCH-136-FLAT-TAIL-ACYCLICITY makes \(\overline K\)
acyclic. The long exact sequence of (136.FT6) now says exactly
that multiplication by \(p\) is bijective on \(H^j(K)\).

If the rational complex is acyclic, every class of \(H^j(K)\)
is \(R\)-torsion and is therefore killed by a power of \(p\).
Injectivity of multiplication by \(p\) forces the class to be
zero. No finite generation or Nakayama argument is needed.
\end{proof}

\paragraph{ARCH-136-FREE-ARITY-TAIL-OBSTRUCTION.}
Let \(R=\mathbb Z_{(p)}\), \(S_R=\Gamma_R(T)\), and let
\(E\) be a nonnegative weight/form graded differential
\(S_R\)-algebra with a further direct-sum grading
\[
                  E=\bigoplus_{r\geq0}E(r).
                                                        \tag{136.FT7}
\]
Assume that the differential preserves this arity, multiplication
adds it, coefficients have arity zero, and the following conditions
hold for each fixed positive \(r\):
\begin{enumerate}
\item the complex \(E(r)\) is bounded in form degree;
\item its terms are weight-graded flat over \(S_R\);
\item each \(H^q(E(r))\) vanishes in all sufficiently high weights.
\end{enumerate}
Then every \(H^q(E(r))\) has bijective multiplication by \(p\).
In particular every native \(p\)-torsion class from a homogeneous
subalgebra must become a boundary in \(E\).

These assumptions hold pointwise for an arity-preserving free
graded-PD envelope having only finitely many form degrees in each
arity and satisfying the archive's vertical admissibility condition.
The assertion also applies when an arity block is replaced by a
bounded complex of flat modules through a weight-preserving
quasi-isomorphism.

\begin{proof}
Only finitely many form degrees occur at arity \(r\).
Take the maximum of the finitely many bounds in condition 3.
All original weights are nonnegative, so the cohomology of
\(E(r)\) lies in one finite weight band. Apply
ARCH-136-INTEGRAL-FLAT-TAIL.

For a free graded-PD algebra, its divided-power and exterior monomials
give an \(S_R\)-basis, and the monomials of fixed arity give a basis
of that arity block. A homogeneous free extension retains this
description. In an admissible algebra, high-weight cohomology of
every positive arity is zero: all positive-form cohomology vanishes,
and the degree-zero coefficient isomorphism has image in arity zero.
This checks the claimed application of the three conditions.
The quasi-isomorphic version follows by applying the integral lemma
to the stated bounded flat representative.
\end{proof}

\paragraph{An actual prescribed-map obstruction.}
Consider the odd unit disk
\[
 D=S_R\otimes_R\Lambda_R(a)\otimes_R\Gamma_R(b),
 \quad da=b,\quad |a|_{\rm form}=1,\quad |b|_{\rm form}=2,
 \quad\operatorname{wt}(a)=\operatorname{wt}(b)=m\geq2,
                                                        \tag{136.FT8}
\]
with both \(a,b\) of auxiliary arity one. In arity \(p\) its
two-term positive complex is
\[
 S_R\,a\gamma_{p-1}(b)
       \xrightarrow{\ p\ }S_R\,\gamma_p(b)
 \quad\text{in forms }2p-1,2p.
                                                        \tag{136.FT9}
\]
Thus \(z=\gamma_p(b)\) is a native class killed by \(p\).
ARCH-136-NATIVE-LIFT supplies an actual admissible theory \(C\)
and a guarded pure horizontal source map whose chosen simplex-one
disk has these data, and for which the image of \(z\) is nonzero
in \(H^{2p}(C_1^{mp,*})\cong\mathbb Z/p\) on its extra summand.
Localizing that construction at \(p\) preserves the nonzero class.

Consequently no factorization of this disk map
\[
                        D\longrightarrow E\longrightarrow C_1
                                                        \tag{136.FT10}
\]
can satisfy the three arity-block conditions above, with the first
map preserving the displayed arity. They would make \(z\) a boundary
in \(E\), contrary to its image in \(C_1\).
This already prevents a pointwise, linearly natural completion of
all disk algebras by such bounded flat arity blocks from retaining
all guarded target maps. Under linear naturality, a selected disk
is also an actual retract of a free collection of disks, so the
same test survives adding other disk generators in that framework.

\paragraph{Implication for a derived radial envelope.}
For finitely many form-two disk generators, every monomial in the
original algebra of total arity \(r\) has form degree at most
\(2r\). The radial image preserves this arity. Resolving its
homogeneous relations by free generators in one lower cochain
degree preserves the same upper bound at each fixed arity:
products add both arity and form degree, and a new primitive has
lower form degree than its specified differential. Retaining a
nonnegative lower bound therefore leaves a bounded flat complex
at each arity. If such a resolution were modified only in sufficiently
high coefficient weights so as to be vertically admissible and to
retain every guarded target map, (136.FT10) would contradict
ARCH-136-FREE-ARITY-TAIL-OBSTRUCTION.

This is the precise force of coefficient-tail syzygies in this
homogeneous free implementation. A growing lower cutoff such as
\(p-1\) does not remove the obstruction. It is not enough to replace
each strict radial relation by a homotopy while retaining a bounded
free complex in its natural arity. A construction may escape the
argument by using nonflat terms, by mixing arities, or by introducing
an unbounded form range within a fixed arity. We have not proved
that any of those alternatives succeeds or fails.

There is also an elementary preliminary compatibility requirement.
If \(\alpha=a_1\cdots a_p\) and one writes
\(dz=py-\alpha\), then \(d^2z=0\) requires
\(p\,dy=d\alpha\). For disk generators the right-hand side is
generally not divisible by \(p\) in the original free algebra.
Thus the entire differential presentation of the radial image,
not one isolated divisibility equation, must be resolved. This
observation alone is not used as an obstruction theorem.

\paragraph{Primary source and actual check scope.}
The characteristic-\(p\) coefficient decomposition (136.FT3) is
Cartan, \emph{Puissances divis\'ees}, S\'eminaire H.~Cartan
1954--1955, Expos\'e 7, \S7, printed p.~7-10, the displayed
truncated-polynomial decomposition following equation (16):
\url{https://www.numdam.org/item/SHC_1954-1955__7_1_A7_0.pdf}.
The primary page and formula were inspected directly in this run;
the decomposition is also proved above. The bounded-flat theorem
uses the explicitly written periodic resolution and double complex,
not a claimed general dg-PD model structure. The full guarded native
example is the already checked ARCH-136-NATIVE-LIFT; its proof was
read in full again for this application. No source theorem is being
claimed for the new arity obstruction.
